\subsection{A common continuation for different economic decisions}
\label{sec:r16menumain}
The capital economy permits a direct study of continuation reuse. Consider a
finite-period decision in which a temporary consumption weight, consumption
externality, or consumption cap changes at the current date. Subsequent preferences,
technology, Gaussian innovations, terminal dispersion cost, and reference
consumption remain fixed. For task $\theta=(w,\eta,u)$, the action-dependent
payoff is
\begin{align}
 Q_\theta(y,a)&=A_0\{w\overline{\log a}-\eta\bar a^2/2\}
                   -W_0\bar a+S\{x_a(y)\},\notag\\
 x_a(y)&=y+h\{c\bm1+f(y)-a\},\qquad
 S(x)=\E F(x;\xi),\qquad a\in[\underline a,u]^d.
 \label{eq:r16mainmenu}
\end{align}
Here a known current term independent of $a$ is suppressed. The random
variable $F(x;\xi)$ includes the first Gaussian innovation at postdecision
mean $x$, all subsequent reference rewards, and the terminal payoff. The
stored period weights and innovation amplitudes define this finite capital
problem. Its decision values have their own interpretation, separate from
the continuous-diffusion value in the preceding payoff theorem.

Proposition~\ref{prop:r16menucache} establishes that $S$ is identical across
these tasks. NBO therefore fits one scalar continuation and queries its
gradient when the current economic problem changes. This reuse has a precise
information boundary. A stored payoff or costate remains an exact rollout
observation at the argument where it was computed. Reusing the same innovation
array at another argument requires reevaluating its state recursion and
payoff. The Raw comparators may cache valid observations and innovations under
the same rule. The comparison concerns the accuracy and complete cost of
these representations of the common continuation.

The scalar representation also imposes structure on costate fitting. Fix
postdecision arguments and positive weights before drawing rollout labels.
Stack the true gradients, multiplied by the square roots of the weights, in
$q$. Let the corresponding labels be $z=q+\varepsilon$, with
$\E\varepsilon=0$ and finite covariance $\Sigma$. For differentiable scalar
features fixed independently of these labels, let $P$ be the Euclidean
orthogonal projection onto their stacked weighted gradient columns. A
least-squares scalar critic has fitted gradients $Pz$, and
\begin{equation}
 \E\|Pz-q\|^2-\E\|z-q\|^2
   =\|(I-P)q\|^2-\operatorname{tr}\{(I-P)\Sigma\}.
 \label{eq:r16mainprojection}
\end{equation}
Proposition~\ref{prop:r16menuprojection} proves the identity without requiring
independent coordinates or Gaussian label noise. The variance removed outside
the gradient dictionary is the statistical benefit; approximation error is
its cost. When the true gradient belongs to the dictionary and the removed
variance is positive, the scalar critic strictly reduces this risk. The same
fitted object can then serve every task in~\eqref{eq:r16mainmenu}.

This calculation gives a conditional statistical account of the scalar
critic. A fully trained nonlinear network generally chooses its features
using its labels and need not coincide with this fixed projection. A Raw
procedure using the same projection has the same identity. The economic
comparisons consequently evaluate all five complete procedures on independent
payoff observations, including their approximation errors, action choices,
and computational work. Technical Section~\ref{sec:r16menu} gives the full
proofs, direct conditional payoff range, and separate accuracy account over
an entire scalar intervention interval.
